\section{Minimal posterior-and-report-predictive coordinates}
\label{sec:v9-saturation}
The variational theorem does not require a sufficient statistic. When
such a statistic is desired, the finite instrument admits a direct
construction. It distinguishes the cost of prediction from the cost of
preserving the whole parameter experiment.

Fix the finite instrument of Section~\ref{sec:v9-minimax}. For an
action--report history $h=(a_0,y_1,\ldots,a_{t-1},y_t)$ define the
unnormalized forward vector by
\begin{align}
 u_{\theta,0}(z)&=\nu_\theta(z),\nonumber\\
 u_{\theta,t+1}(hay,z')
   &=\sum_z u_{\theta,t}(h,z)K_{\theta,t}(z',y\mid z,a),
   \qquad L_\theta(h)=\sum_z u_{\theta,t}(h,z).
 \label{eq:v9-forward}
\end{align}
Choose a reference prior $\pi_\theta>0$. A history is supported if
$\overline L(h):=\sum_\theta\pi_\theta L_\theta(h)>0$. On supported
histories put
\begin{align}
 \lambda_\theta(h)&=\pi_\theta L_\theta(h)/\overline L(h),\nonumber\\
 k_t(y\mid h,a)&=\overline L(hay)/\overline L(h).
 \label{eq:v9-finitepredict}
\end{align}
The likelihoods in these formulae do not contain action-selection
probabilities. Under any parameter-independent observation-based policy,
those probabilities cancel from the posterior and from likelihood ratios.
Thus the formulae apply to every policy on its supported histories, not
only to a policy used to explore the history tree.

A history coordinate $s_t(h)$ is called \emph{saturated and predictive}
if it determines $\lambda(h)$, determines every vector
$k_t(\cdot\mid h,a)$, and has a deterministic update from $(s_t,a,y)$ on
supported continuations. At time $T$ only the first condition is imposed.
This is the finite, terminal-parameter version of the three conditions
in Section~\ref{sec:v8-sufficiency}. Additional latent-target marks
require the corresponding conditional kernels as further decorations;
parameter sufficiency alone does not imply their sufficiency.
Sections~\ref{sec:v10-marked}--\ref{sec:v10-synthesis} supply the
marked reconstruction theorem and an output-sensitive construction.

Construct partitions backwards. At level $T$ group histories with the
same vector $\lambda(h)$. Having constructed the class map $c_{t+1}$,
let $c_t(h)$ be the class determined by the finite tuple
\begin{equation}
 \left(\lambda(h),\,
   \big(k_t(y\mid h,a),\,c_{t+1}(hay)\big)_{a,y}\right).
 \label{eq:v9-saturatedkey}
\end{equation}
For an unsupported continuation its class entry is one fixed cemetery
symbol. Let $K_t$ be the number of classes at level $t$.

\begin{theorem}[Minimal posterior-and-report-predictive coordinate]
\label{thm:v9-saturation}
For every finite controlled instrument the partitions
\eqref{eq:v9-saturatedkey} form a saturated predictive coordinate. Every
other such deterministic history coordinate refines them at each level.
Consequently $K_t$ is its least possible width, simultaneously at every
level. The partitions, as equivalence relations on supported histories,
are independent of the choice of the strictly positive reference prior.

There is a parameter-independent finite reconstruction kernel
$L_t(y\mid c,a,c')$ which, from successive quotient classes and the
chosen controls, reconstructs the report stream under every admissible
policy. The original and quotient parameterized report experiments have
mutually exact causal simulations. This assertion concerns reports and
parameter likelihoods only; latent-target preservation is the separate
marked statement of Theorem~\ref{thm:v10-reconstruction}. Their additional persistent state is at most
$\max_tK_t$ labels, so composing either simulation with an $S$-state
encoder costs at most $(\max_tK_t)S$ states.

If the instrument and prior have rational entries, the classes and
reconstruction kernels can be computed by exact finite arithmetic.
The calculation is polynomial in the size of the \emph{expanded finite
history tree} and its rational input length. That tree, and the least
width $K_t$, may grow exponentially with the horizon.
\end{theorem}
\begin{proof}
Equality of the keys at level $t$ implies equality of the posterior
vectors, of all reference predictive probabilities, and of the class of
every common positive-probability successor. Sending $(c,a,y)$ to that
successor class therefore defines an update. Unsupported continuations
may be assigned arbitrarily. This proves the three required properties.

For minimality, at time $T$ any saturated coordinate separates unequal
posterior vectors, so its fibres refine $c_T$. Suppose the assertion has
been proved at $t+1$. Equal coordinates at time $t$ have equal posterior
and predictive vectors. Their supported successors under any fixed
$(a,y)$ have equal next coordinates by recursion, and hence equal
$c_{t+1}$ by the induction hypothesis. Their complete keys
\eqref{eq:v9-saturatedkey} coincide. Backward induction proves the
claimed simultaneous refinement, rather than just a bound on a terminal
statistic.

To check prior independence, first observe that equality of posterior
vectors under a positive prior is equivalent to proportionality of the
nonnegative likelihood vectors $(L_\theta(h))_\theta$. This relation is
independent of the prior. For a fixed parameter with
$\lambda_\theta(h)>0$, the actual next-report probabilities are
\begin{equation}
 P_\theta(y\mid h,a)
   =k_t(y\mid h,a)\frac{\lambda_\theta(hay)}{\lambda_\theta(h)}.
 \label{eq:v9-finitechangeofmeasure}
\end{equation}
In an existing class the right side is constant on histories for every
supported successor, because the successor class determines its
posterior. The set of parameters with positive likelihood is also
constant on the class. Under a different positive prior, its new
posterior weights are functions of its old posterior weights; averaging
these common parameter-specific predictive kernels therefore gives a
common new reference predictive kernel. The existing partition is
saturated, predictive and recursive for the new prior. Minimality in
both directions proves equality of the two partitions.

For reconstruction, write $F_t(c,a,y)$ for the class update and define
\begin{align*}
 H_t^0(c'\mid c,a)&=\sum_{y:F_t(c,a,y)=c'}k_t(y\mid c,a),\\
 L_t(y\mid c,a,c')&=
 \frac{\1_{\{F_t(c,a,y)=c'\}}k_t(y\mid c,a)}{H_t^0(c'\mid c,a)}
                 \quad\text{if }H_t^0(c'\mid c,a)>0.
\end{align*}
Choose an arbitrary probability vector when the denominator is zero.
The posterior at the next class is constant over the summed reports.
Equation~\eqref{eq:v9-finitechangeofmeasure} thus factors the actual
joint law of report and next class into this same $L_t$ and the
parameter-specific class-transition probability. Conditioning further
on the next class cancels the likelihood factor. This gives a single
parameter-independent report reconstruction, including when parameters
have different supports. Histories having zero likelihood under a
parameter are not visited under that parameter.

Starting from the initial class, the reconstruction stores the previous
class, receives the next class after the actual chosen action, samples
its report using $L_t$, and replaces its stored class. A target policy
then chooses its next action from the reconstructed reports. Induction
in time proves equality of the resulting action--report laws for every
parameter and policy. In the other direction, apply $F_t$ to the original
reports. Product registers give the displayed storage bound; neither
simulation presumes a cost-free transcript.

Finally, \eqref{eq:v9-forward} enumerates the rational forward vectors
on a finite tree. Formulae \eqref{eq:v9-finitepredict} and
\eqref{eq:v9-saturatedkey} use only finite sums, products, divisions by
positive numbers, and equality tests of rational tuples. Class labels
can be assigned by sorting those tuples backwards. To see the stated
bit-complexity bound, take a common denominator for all rational input
entries. Its bit length is at most their total input length. At level
$t$, a forward entry can be represented with the $t$th power of this
denominator times the initial denominator; sums add at most the bit
length of the number of summands at each step. Thus every intermediate
integer has bit length polynomial in horizon and input length.
Posterior divisions and tuple comparisons preserve a polynomial bound.
There are only as many forward and backward steps as nodes and edges
of the expanded tree, times its finite hidden and parameter dimensions.
This proves effectivity without any oracle for exact real comparisons.
\end{proof}

The theorem constructs the object assumed by the causal factorization
in this finite class. Its minimality is among the \emph{specified}
likelihood-saturated, raw-report-predictive coordinates. It is not a claim
that every exact simulator, or every task-specific lossy encoder, must
use these coordinates. The variational principle optimizes the latter
directly and may use substantially less memory.

\begin{example}[Saturation can have an exponential price]
\label{ex:v9-saturationcost}
Let the parameter range over $\{0,1\}^d$, and reveal it noiselessly one
bit at a time during the first $d$ observations. Thereafter all reports
are independent fair bits. At time $d$ all future-report predictions
coincide, so the future-report predictive quotient has one class.
Nevertheless the posterior on the parameter is a different point mass
for each word. Every saturated coordinate at that time has at least
$2^d$ classes, and the construction above has exactly $2^d$.
Predictive closure alone therefore supplies no controlled-cost
parameter saturation. This is a cost calculation in a concrete model,
not an assumption that the missing saturation has bounded size.
\end{example}

\begin{example}[Likelihoods alone need not close prediction]
\label{ex:v9-predictioncost}
There is one parameter. Let $Y_1$ be a fair bit and put $Y_2=Y_1$.
The constant coordinate is recursive and contains every parameter
likelihood, but its one-step predictive law at time $1$ is not a function
of that coordinate. The least saturated predictive widths are
$(K_0,K_1,K_2)=(1,2,1)$. The last quotient class can discard the bit,
because reconstruction of $Y_2$ still uses the previous class. This
example separates the hypotheses of the factorization theorem; it does
not assert impossibility of other simulators with other charged resources.
\end{example}
